\section{Authorizing the Field That Returns Anything}
\label{sec:ch15-the-field-that-returns-anything}
\index{node@\code{Query.node}!authorizing it|(}%

\code{Query.ordersByCustomer} is guarded now. It carries a scope the router
checks, a policy the subgraph checks, and an ownership test in the resolver,
because a scope can say that you may read orders and cannot say whose. Ask for
somebody else's history with a perfectly good token and the graph says no:

\begin{minted}[fontsize=\scriptsize,breaklines]{json}
{"errors":[{"message":"Failed to fetch from Subgraph 'ordering'.","extensions":
  {"errors":[{"message":"You may only read your own orders.","path":["ordersByCustomer"],
   "extensions":{"code":"AUTH_NOT_AUTHORIZED"}}],"serviceName":"ordering",
   "statusCode":200}}],"data":null}
\end{minted}

Here is that same graph, one second later, with the same token, before the fix
this section is about. Nothing below reproduces at tag \code{ch15}, and that is
the point of printing it:

\begin{minted}[fontsize=\scriptsize,breaklines]{json}
{"data":{"node":{"__typename":"Order","id":"T3JkZXI6AAAA0AAAAECAAAAAAAAAAw==",
  "total":{"amount":356.00,"currency":"EUR"},"placedAt":"2026-02-18T11:05:00Z"}}}
\end{minted}

\index{node@\code{Query.node}!the second door}%
Chapter~\ref{ch:hard-modeling-problems} gave this graph its \code{node} field
back. What it also did, and nobody noticed for two chapters, was open a second
door to every addressable row in the graph, and that door does not pass the room
the first one is guarded in. Figure~\ref{fig:ch15-two-doors} is the two paths.

\begin{figure}[htbp]
  \centering
  % One row, reached two ways. Referenced from chapter 15. Bare tikzpicture;
% the figure environment, caption and label live at the call site.
%
% The point of the picture is that guarding the left door does nothing for the
% right one. The right door goes through a service that owns no data, which can
% refuse a type and can never refuse a row, and the router then fetches the row
% itself through _entities - a path that never goes near ordersByCustomer or
% its checks. The dashed box is the checkpoint that did not exist until the gap
% was measured: before it, the right door handed over what the left door had
% refused the same caller a second earlier.
\begin{tikzpicture}[
    font=\small,
    box/.style={draw, rounded corners=2pt, minimum width=22mm,
                minimum height=7mm, align=center, font=\scriptsize},
    gate/.style={draw, rounded corners=2pt, fill=black!8, minimum width=24mm,
                 minimum height=8mm, align=center, font=\scriptsize},
    fixed/.style={gate, dashed},
    flow/.style={-{Stealth[length=2mm]}},
    note/.style={font=\scriptsize, gray, align=center, inner sep=1pt}
  ]

  \node[box] (client) at (0.8,0) {client};
  \node[box] (router) at (5.4,0) {router};
  \draw[flow] (client) -- (router);

  % -- left door: a root field on the service that owns the data -------------
  \node[box] (fieldL) at (2.2,-1.8) {\code{ordering}\\\code{ordersByCustomer}};
  \node[gate] (gateL) at (2.2,-3.7) {a scope,\\then \code{CustomerId}\\\code{== the subject}};

  \draw[flow] (router) -- (fieldL);
  \draw[flow] (fieldL) -- (gateL);

  % -- right door: a stub from a service with no data ------------------------
  \node[box] (fieldR) at (8.8,-1.8) {\code{nodes}\\\code{node(id:)}};
  \node[gate] (gateR) at (8.8,-3.7) {the type only,\\decoded from\\the identifier};
  \node[fixed] (gateE) at (8.8,-5.6) {\code{\_entities}, then\\\code{CustomerId}\\\code{== the subject}};

  \draw[flow] (router) -- (fieldR);
  \draw[flow] (fieldR) -- (gateR);
  \draw[flow] (gateR) -- (gateE);

  \node[note, anchor=west] at (10.4,-3.7) {no data here,\\so no row\\to check};
  \node[note, anchor=west] at (10.4,-5.6) {\code{ordering} again,\\reached by\\the router};

  % -- one row ---------------------------------------------------------------
  \node[box] (row) at (5.4,-7.3) {the \code{Order} row};
  \draw[flow] (gateL) -- (row);
  \draw[flow] (gateE) -- (row);

  \node[note, anchor=north, align=center] at (5.4,-7.9)
    {the dashed check was added last, and until it existed this door\\
     handed over what the other one had just refused};

\end{tikzpicture}

  \caption{Two routes to one order. On the left a root field on the service that
  owns it, where a scope and an ownership check both apply. On the right
  \code{Query.node}, which is answered by a service with no database: it decodes
  the identifier, hands back a stub, and the router fetches the rest through
  \code{\_entities}, a path that has never heard of \code{ordersByCustomer} or
  its rules. Guarding one door does nothing for the other.}
  \label{fig:ch15-two-doors}
\end{figure}

\subsection*{The service that owns nothing cannot check very much}

\index{Mosaic.Nodes@\code{Mosaic.Nodes}!what it can and cannot decide}%
The first place to look for a fix is where the field lives, so look at what is
there. Three of the four node resolvers in \code{Mosaic.Nodes} are two lines,
and the one for orders was too, until this chapter. It was this:

\begin{minted}[fontsize=\footnotesize,breaklines]{csharp}
[NodeResolver]
public static Order ResolveOrder(Guid id) => new() { Id = id };
\end{minted}

No database, no domain, no \code{\_entities}. Decision 75 in this book's spec
made that a feature: a service that owns nothing is the honest home for a field
that has to be able to return anything. It is also a service that cannot answer
\enquote{is this order yours}, and never will, because it has never seen an
order.

What it can answer is what \emph{kind} of thing an identifier names, because
decoding that is the entirety of its job and the type name is inside the
identifier. So it gets the only rule it is capable of holding:

\begin{minted}[fontsize=\footnotesize,breaklines]{csharp}
[NodeResolver]
public static Order? ResolveOrder(Guid id, ClaimsPrincipal? user)
    => user?.Identity?.IsAuthenticated == true
        ? new Order { Id = id }
        : null;
\end{minted}

\index{ClaimsPrincipal@\code{ClaimsPrincipal}!nullable or throwing}%
That is coarse and it is deliberate. Orders are not public, and this is the last
point at which anybody can say so cheaply. Note the question mark on the
parameter, which is not decoration: HotChocolate's parameter binder hands a
nullable \code{ClaimsPrincipal} null for an anonymous caller and throws an
\code{ArgumentException} out of a non-nullable one, so the difference between
the two spellings is the difference between \enquote{nobody is signed in} and an
execution error.

\index{node@\code{Query.node}!what the identifier alone answers}%
It matters more than it looks, and the measurement that convinced me is this
one. Before that check existed, an \emph{anonymous} request could ask
\code{node(id:)} for somebody's order and get an answer, as long as it selected
only \code{id} and \code{\_\_typename}. No subgraph was called at all. Both of
those fields are carried by the argument itself, so the stub was the whole
answer, and every guard in Ordering was irrelevant because Ordering was never
asked. A graph can leak the existence and type of a row without any service
touching a database.

\subsection*{Where the instance check has to go}

\index{reference resolvers!an ownership check inside one}%
Whose order it is stays with Ordering, in the reference resolver the router
reaches through \code{\_entities}:

\begin{minted}[fontsize=\footnotesize,breaklines]{csharp}
var order = await orderById.LoadAsync(orderId, cancellationToken);

if (order is null || !OrderAccess.IsSelf(user, context, order.CustomerId))
{
    return null;
}

return order;
\end{minted}

Null rather than an error, and unlike the query field beside it that is the
right answer. An order you may not have and an order that does not exist should
be indistinguishable, or a client with a list of guessed identifiers learns
which of its guesses were real from the difference in the replies. The
specification's \code{[\_Entity]} is nullable for the not-found case and this
borrows it.

Whether that actually holds is a measurement rather than a hope, because
non-null propagation has undone more careful decisions than this one. It holds:
asking for an order that belongs to somebody else and asking for an order that
was never created produce byte-identical responses, and both verification
scripts assert it by comparing the two.

\index{authorization!type against instance}%
So the two checks are not redundant and neither can be moved. The node service
refuses a \emph{type} to a stranger and could not refuse an instance if it
wanted to. Ordering refuses an \emph{instance} and never sees the request that
the node service turned away. Each rule ended up in the only service that holds
what it needs to ask about.

\index{header forwarding!what the reference resolver depends on}%
And the reference resolver's check depends on something invisible from the file
it is written in. It can only work if the caller's token reached that process,
which is a router configuration two directories away. Take the \code{headers}
block out and every line of it still runs, sees nobody, and refuses the owner
their own order. That is a failure mode with no error message anywhere in it,
and the only reason I know what it looks like is that I built the check before I
built the configuration.

\medskip
Three layers, four fields, and two doors closed. What is left is the case where
none of it works, and it is the one the previous chapter handed over.

\index{node@\code{Query.node}!authorizing it|)}%
